\section{Objetivos}

\begin{itemize}
	\item Comprender y aplicar las técnicas de escalamiento de magnitud y de frecuencia en circuitos eléctricos para modificar los valores de sus componentes y su respuesta operativa según requerimientos de diseño específicos.
	\item Aplicar el factor de escalamiento de magnitud $K_m$ para ajustar los componentes de una red a valores más prácticos o comerciales, garantizando que la respuesta en frecuencia permanezca inalterada.
	\item Implementar el escalamiento de frecuencia mediante el factor $K_f$ para desplazar la respuesta de un circuito (como frecuencias de corte o resonancia) a un punto deseado del espectro.
	\item Determinar analíticamente las funciones de transferencia de los circuitos prototipo y escalados mediante el análisis de mallas para verificar su comportamiento matemático.
	\item Validar, a través de herramientas de simulación, que las relaciones de impedancia de entrada y los diagramas de Bode correspondan con los factores de escala calculados teóricamente.
	\item Analizar el efecto de los factores de escalamiento sobre los elementos reactivos del circuito, observando cómo varían los valores de inductancia y capacitancia para mantener la respuesta del sistema.
\end{itemize}